% ---------------------------------------------------------------------------
% McpRisk v5 -- policy-grade static severity scoring for MCP servers.
% AAMAS 2027 format (26th ICAAMAS, 3-7 May 2027, JW Marriott Hanoi, Vietnam).
%
% AAMAS uses the ACM `acmart' class in sigconf mode with the IFAAMAS copyright
% block. Section numbering is arabic by default here (1, 1.1, 1.1.1) -- no
% override needed, unlike IEEEtran.
%
% Overleaf: upload this folder, compiler pdfLaTeX, main document main.tex.
% `acmart' and ACM-Reference-Format.bst ship with Overleaf's TeX Live.
%
% SUBMISSION vs CAMERA-READY -----------------------------------------------
% AAMAS main-track review is double-blind. For submission, switch the class
% option line below to the `anonymous' variant and keep \settopmatter{
% printfolios=true} so reviewers get page numbers. Confirm the page limit and
% anonymity policy against the AAMAS 2027 submission instructions when they
% are published (deadline listed as early Oct 2026, TBC).
% ---------------------------------------------------------------------------

% -- SUBMISSION (double-blind, page numbers on):
% \documentclass[sigconf, anonymous, review]{acmart}
% \settopmatter{printacmref=false, printfolios=true}

% -- CAMERA-READY / working draft (current):
\documentclass[sigconf]{acmart}
\settopmatter{printacmref=false}

% AAMAS boilerplate: no ACM reference block, no copyright footnote, no running
% headers. These three lines are prescribed by the IFAAMAS template.
\renewcommand\footnotetextcopyrightpermission[1]{}
\pagestyle{plain}

% NOTE ON \setcopyright: the AAMAS template prescribes \setcopyright{ifaamas},
% but `ifaamas' is NOT a value in the stock acmart.cls -- it only exists in the
% patched acmart.cls that IFAAMAS ships inside the official AAMAS template
% bundle. Compiling against stock acmart with it raises
%   "Package xkeyval Error: value `ifaamas' is not allowed."
% So this draft uses a stock-safe value. When you drop the official AAMAS
% bundle's acmart.cls into this folder, swap the two lines below.
% \setcopyright{ifaamas}
\setcopyright{rightsretained}

\acmConference[AAMAS '27]{Proc.\ of the 26th International Conference on
Autonomous Agents and Multiagent Systems (AAMAS 2027)}{May 3 -- 7,
2027}{Hanoi, Vietnam}
\acmYear{2027}
\copyrightyear{2027}
\acmDOI{}
\acmPrice{}
\acmISBN{}
\acmSubmissionID{}

% --------------------------------------------------------------- packages ---
% acmart already loads amsmath, booktabs, graphicx and xcolor -- loading
% amssymb on top of it clashes with \Bbbk from the Libertine math setup.
\usepackage{balance}
\usepackage{enumitem}   % optional args on enumerate (contributions list)

% ---------------------------------------------------------------- shortcuts ---
\newcommand{\eg}{e.g.,\ }
\newcommand{\ie}{i.e.,\ }
\newcommand{\framework}{\textsc{McpRisk}}

% Placeholder marker for sections not yet written. Delete the macro once the
% skeleton is filled in; it renders visibly so nothing ships by accident.
\newcommand{\stub}[1]{\textit{\textcolor{gray}{[To be written: #1]}}}

\begin{document}

% ------------------------------------------------------------------- title ---
\title{Deriving Design-Time Risk Severity for Model Context Protocol Servers
from Organizational Policy Alone}

\author{Tomer Ovadya}
\affiliation{%
  \institution{Ben-Gurion University of the Negev}
  \department{Software and Information Systems Engineering}
  \city{Beer-Sheva}
  \country{Israel}}
\email{ovadyat@post.bgu.ac.il}

\author{Asaf Shabtai}
\affiliation{%
  \institution{Ben-Gurion University of the Negev}
  \department{Software and Information Systems Engineering}
  \city{Beer-Sheva}
  \country{Israel}}
\email{shabtaia@bgu.ac.il}

\renewcommand{\shortauthors}{Ovadya and Shabtai}

% ---------------------------------------------------------------- abstract ---
\begin{abstract}
\stub{abstract --- write last, once the evaluation section is fixed.}
\end{abstract}

% ------------------------------------------------------------ CCS + keywords ---
\begin{CCSXML}
<ccs2012>
   <concept>
       <concept_id>10002978.10003006.10003013</concept_id>
       <concept_desc>Security and privacy~Distributed systems security</concept_desc>
       <concept_significance>500</concept_significance>
       </concept>
   <concept>
       <concept_id>10002978.10002991.10002995</concept_id>
       <concept_desc>Security and privacy~Access control</concept_desc>
       <concept_significance>300</concept_significance>
       </concept>
   <concept>
       <concept_id>10010147.10010178.10010179</concept_id>
       <concept_desc>Computing methodologies~Intelligent agents</concept_desc>
       <concept_significance>300</concept_significance>
       </concept>
 </ccs2012>
\end{CCSXML}

\ccsdesc[500]{Security and privacy~Distributed systems security}
\ccsdesc[300]{Security and privacy~Access control}
\ccsdesc[300]{Computing methodologies~Intelligent agents}

\keywords{Model Context Protocol; AI agents; risk scoring; security
categorization; data classification policy; LLM security; access control}

\maketitle

% ---------------------------------------------------------------------------
% Section 1 -- Introduction (short paper: FLAT, no subsections)
%
% Structure follows the "technical-challenge version 2" template from the
% research-paper-writing skill: our insight has historical roots in a
% traditional line (risk-adaptive access control, 2007-2011), so that line is
% used as conceptual backing before showing why modern MCP defenses still fail.
%
% Paragraph roles:
%   P1 opening    -- agentic applications and MCP
%   P2 direction  -- server as protected resource; misuse through legitimate access
%   P3 challenge  -- RBAC/ABAC answer "may it?", not "how dangerous is it?"
%   P4 challenge  -- risk-based access control already solves this shape, but
%                    not for agents; current MCP defenses remain binary
%   P5 insight    -- what the risk must be computed from (static, dynamic, history)
%   P6 method     -- our methodology
%   P7 evidence   -- implementation, evaluation, three contributions
%
% Budget: ~850 words. Citations: 24 of the 37 in refs.bib.
%
% EXPECTED: \ref to sec:framework / sec:eval typesets as "??" until those
% sections exist. Not a build error.
% ------------------------------------------------------------------------
\section{Introduction}
\label{sec:intro}

AI agents no longer only generate text; they act. They read data, modify files,
update databases, and call external services for their users. The Model Context
Protocol (MCP)~\cite{anthropic2024mcp} standardizes this exchange: a server
publishes tools, and a connected agent discovers and invokes them over
JSON-RPC~\cite{mcp2026spec}. The public server ecosystem is growing quickly~\cite{hou2026mcp}, and audits of deployed servers find recurring weaknesses~\cite{hasan2025firstglance}. An agent's decision
now lands directly on the systems and resources behind those tools. These tools let agent decisions directly affect organizational resources. MCP
standardizes how tools are discovered and invoked, but it does not assess the
risk of a specific call. Its annotations describe the tool rather than the
resource it affects, so a destructive tool receives the same label whether it
targets a temporary file or a production database.

Most MCP and agent security research studies how malicious servers or tools can
harm the agents that use
them~\cite{invariant2025toolpoisoning,shi2025toolhijacker}. Other work audits
MCP server implementations and exposed capabilities for software weaknesses or
excessive privilege~\cite{kumar2026mcpinsos,huang2026auditing}. At runtime,
MTGuard inspects invocation parameters and observed execution to detect policy
violations~\cite{he2026mtguard}. These approaches identify vulnerabilities,
excessive permissions, or unsafe behavior. They do not estimate the consequence
of an authorized invocation from the sensitivity and scope of the resource it
affects. Our work provides this estimate.


We study \emph{misuse}: a legitimate capability used beyond what the task requires~\cite{anderson1980threat}. It needs no malicious agent and no malicious server. It can start with an operator who states a task loosely, delegates too broadly, or grants excess permission~ \cite{owasp2025llm06} . Studies of MCP
applications and servers report broad system access and over-privileged
tools~\cite{li2025privilege,huang2026auditing}. Misuse can also result from
model behavior: an agent may select an available but inappropriate
tool~\cite{yin2025reasoningtrap}, or recognize a risk in general but fail to
account for it during execution~\cite{tang2025riskknowledge}. We do not ask why the agent made the call. We assess what the call can do.
Reading a build log and reading stored credentials may use the same tool, but
the risks are different. We therefore analyze the full invocation: the tool,
the resource it affects, the arguments, the execution context, and the agent's
recent actions.



Access control decides whether a request is allowed. Role and attribute-based
models~\cite{sandhu1996rbac,hu2014abac} evaluate permissions under a policy,
but they do not normally assign a risk level to an allowed request. This gap
also appears in agent frameworks: exposing a tool does not ensure that each call
is authorized using its actual argument values~\cite{mellafe2026capability}.
Authentication identifies the caller, and authorization defines what the caller
may do. Neither measures the possible consequence of a specific authorized
call.


Adding risk to an access-control decision is not new. Early work quantified risk
during authorization~\cite{cheng2007fuzzymls}, and later work incorporated
risk-adaptive access control into an attribute-based
framework~\cite{kandala2011radac}. The area has since been widely
studied~\cite{atlam2020riskbased}. Recent work extends this idea to agents,
using an LLM judge to consider task risk, resource risk, and model
uncertainty~\cite{fleming2025tbac}. However, these models do not define an
MCP-specific consequence score for a concrete invocation based on the tool, the
sensitivity of the affected resource, and the scope of the call.


To our knowledge, no existing method estimates the consequence of a concrete
MCP invocation from the sensitivity and scope of the resource it affects. Many
defenses operate at the agent, client, or middleware layer, where they restrict
available tools or authorize calls before
execution~\cite{shi2025progent,zhu2025miniscope,betser2026agentrim,
li2026conleash,abaev2026agentguardian}. Server-side authorization has also begun
to appear: IGAC checks a proposed call and its payload against the user's
declared intent~\cite{zhu2026igac}. Other runtime systems return safety scores
or decisions such as allow, review, or
block~\cite{yang2026agenttrust,zhang2026stars,hossain2026nexus}. These outputs
indicate whether a call is permitted, suspicious, or requires intervention, but
they do not estimate the severity of an authorized call for the specific
organizational resource it affects.

Vulnerability-scoring frameworks do not address this problem either. CVSS
scores the severity of software, hardware, and firmware
vulnerabilities~\cite{first2023cvss}, while AIVSS extends vulnerability scoring
with factors specific to agentic AI systems~\cite{huang2026aivss}. Neither
estimates the consequence of a concrete MCP invocation. Our method provides
this consequence score and uses it to allow, review, or refuse each call.


Risk belongs to a tool--resource pair, not to the tool alone. For example,: \texttt{read\_file} poses little risk when applied to a public document but
substantial risk when applied to regulated data. We therefore compute a
baseline from three factors: tool impact, resource sensitivity, and
\emph{blast radius}, meaning the portion of the resource affected by one call.
The call's arguments may change this baseline, while its context and the
agent's recent activity may indicate whether the call is unusual.

We compute risk at two stages. At design time, we assign a baseline score to
each tool--resource pair exposed by the server. The score combines tool impact,
resource sensitivity, and blast radius. Resource sensitivity is derived from
the organization's information-classification policy, following established
guidance for mapping information types to impact
categories~\cite{nist2004fips199,nist2008sp80060}. At request time, we update
the baseline using the call's arguments, execution context, and the agent's
recent activity. The result is an auditable score an organization can compare
with thresholds of its own choosing; deciding whether to approve a call,
require human review, or refuse it remains the organization's
(\S\ref{sec:framework}).

Adjustable autonomy studies when an agent should act independently and when
control should pass to a person~\cite{pynadath2002adjustable}. Our score
provides a measurable basis for this decision. We implement the method as an
enforcement framework and evaluate it using tool catalogs from real vendor MCP
servers, together with organizational policies, resource inventories, and
invocation scenarios (\S\ref{sec:eval}).

This paper makes three contributions. First, a \textbf{risk-analysis methodology}
for assessing the misuse risk of an MCP invocation from tool impact, resource
sensitivity, blast radius, runtime arguments, execution context, and the agent's
recent activity, deriving resource sensitivity from organizational policy rather
than from a fully labeled asset inventory. Second, an \textbf{enforcement
framework} that combines the design-time and request-time assessments and uses
them to make risk-based access-control decisions for MCP servers. Third, an
\textbf{evaluation setting} that pairs tool catalogs from real vendor MCP servers
with an organizational policy, resource registers, and live invocation traffic,
together with a descriptive report of what the framework produces over sixteen
server deployments and an explicit account of what that report does not
establish.
% ---------------------------------------------------------------------------
% Section 2 -- Related Work (short paper: FLAT, four thematic strands)
%
% Not a paper list. Four strands, each one paragraph, each closing on the
% sentence that says why that strand does not answer our question. Bold run-in
% headings instead of \subsection -- subsections cost vertical space and invite
% padding (WRITING-GUIDE.md sec. 7).
%
% Paragraph roles:
%   P1 (a) MCP security        -- the field protects the agent from the server
%   P2 (b) agent/MCP defenses  -- they end at an enforcement verdict
%   P3 (c) access control      -- KEY STRAND: RBAC/ABAC permit; risk-adaptive
%                                 access control conditions on a computed risk
%                                 value. fleming2025tbac is the NEAREST work and
%                                 is named as such -- the gap is the *input*
%                                 (a decomposed per-invocation estimate), not
%                                 the absence of any risk value for agents.
%   P4 (d) risk scoring        -- scores the server, the tool, or the log
%
% Deliberately NOT repeated here (sec. 1 and sec. 3 already carry it, with the
% same keys): the four-route misuse taxonomy -- yin2025reasoningtrap,
% tang2025riskknowledge, zhan2024injecagent, lynch2025misalignment. P2's job is
% what the defenses *produce*, not why they are needed.
%
% Budget: 600-700 words, ~31 citation keys, all from refs.bib.
% ---------------------------------------------------------------------------
\section{Related Work}
\label{sec:related}

\textbf{MCP Server Security and Risk Assessment.} 
Prior work surveys security risks across the MCP
ecosystem~\cite{hou2026mcp} and audits open-source MCP servers for
software weaknesses, MCP-specific vulnerabilities, and maintainability
problems~\cite{hasan2025firstglance}. Several studies examine attacks in which
a malicious server, tool description, or external input influences the
agent~\cite{invariant2025toolpoisoning,shi2025toolhijacker,
zhao2025parasites}. We study the opposite direction: an agent sends a call that
may harm the MCP server or an organizational resource reached through it.
Server-to-agent attacks are outside our scope. We assume only that each tool
performs the operation represented by its declared interface and resource
mapping.

Other work examines the capabilities and weaknesses of MCP servers. Li et al.
analyze 2,562 MCP applications and measure their use of file, network, system,
and memory APIs~\cite{li2025privilege}. Huang et al. combine static pattern
matching with sandboxed execution to detect capabilities such as file access,
network requests, and command execution~\cite{huang2026auditing}. MCP-in-SoS
scans 222 server repositories for code weaknesses, maps them to CWE and CAPEC
entries, and derives repository-level risk from the likelihood and severity
those catalogs assign~\cite{kumar2026mcpinsos}. All three assess the server as
an artifact, independently of the resources a given installation holds, so the
same server receives the same assessment wherever it is deployed.

ConLeash checks whether an MCP call stays within an approved consent boundary~\cite{li2026conleash}. However, it does not estimate the risk or possible impact of a call that is already allowed. Our work addresses this gap by assessing the risk of each invocation from the server side.




\textbf{Defenses for MCP and tool-using agents.} 
IGAC checks whether an MCP invocation and its arguments remain within the user’s stated intent, while MTGuard compares invocation parameters, runtime behavior, and results against safety policies~\cite{zhu2026igac,he2026mtguard}. Progent and MiniScope enforce least-privilege authorization, whereas AgentGuardian learns admissible tool inputs and invocation sequences from normal executions~\cite{shi2025progent,zhu2025miniscope,abaev2026agentguardian}. These systems check whether a tool call follows the user’s request, the security policy, or the agent’s normal behavior. However, a call can pass all these checks and still cause serious damage. For example, an agent may be allowed to delete files, but deleting a temporary file and deleting a critical production file do not have the same risk; our work measures this difference at the MCP server.

\textbf{Access control.} 
RBAC assigns permissions through roles, while ABAC evaluates attributes of the
caller, resource, requested action, and environment
conditions~\cite{sandhu1996rbac,hu2014abac}. These models can express
fine-grained policies, but they do not directly estimate the possible
consequence of a permitted request.

Risk-adaptive access control adds a calculated risk value to the authorization
decision~\cite{cheng2007fuzzymls}. Kandala et al. place risk within an
attribute-based framework, but leave the concrete risk-calculation method
outside their scope~\cite{kandala2011radac}. Ni et al. use fuzzy rules over
subject and object security information to estimate the risk of an access
operation~\cite{ni2010fuzzy}. Shaikh et al. dynamically calculate trust and
risk for each subject--object pair from the user's past
behavior~\cite{shaikh2012dynamic}. Atlam et al. survey the wider family of
risk-based access-control models~\cite{atlam2020riskbased}.

For agentic systems, Fleming et al. use an LLM judge to create a task policy,
combine static risk values for its tools and resources, and estimate
uncertainty~\cite{fleming2025tbac}. A limitation is that risk is computed for
the whole task and is not updated for each call. For example, reading one
database record and exporting a full table may use the same tool risk value.
Our work instead scores each MCP invocation using its runtime details.

\textbf{Risk scoring for MCP and agents.} 
CVSS scores the severity of software vulnerabilities. A study of 196 evaluators
applying CVSS v3.1 found differences between evaluators, and a follow-up found
that 68% assigned a different severity level to vulnerabilities they had
scored before~\cite{first2023cvss,wunder2024cvss}. AIVSS extends vulnerability
scoring to agentic AI systems~\cite{owasp2026aivss}. These frameworks score
vulnerabilities rather than concrete tool calls.

FIPS 199 and NIST SP 800-60 classify information types and systems by the
potential impact of losing confidentiality, integrity, or
availability~\cite{nist2004fips199,nist2008sp80060}. These classifications
provide an organizational view of consequence that our method uses when
assessing the resource affected by a call.

Work closer to our setting studies risk at different stages. MCP-in-SoS scores
weaknesses in an MCP server repository~\cite{kumar2026mcpinsos}, while
AgenTRIM assigns risk levels to tools and applies stronger checks to riskier
ones~\cite{betser2026agentrim}. STARS scores proposed skill invocations using
the request, capability, and runtime context before
execution~\cite{zhang2026stars}. Owiredu-Ashley instead grades the harm of
actions recorded after execution~\cite{owireduashley2026severity}. Together,
These methods answer different risk questions: they assess server weaknesses,
tool capabilities, invocation safety, or harm after execution. They do not
provide a common measure of the consequence to the organization affected by
the call. For example, deleting a temporary cache and deleting payroll records
may use the same tool and both be irreversible, but their organizational impact
is very different. Our work connects each MCP call to the organization’s
resource policy and estimates its risk before execution.

% ---------------------------------------------------------------------------
% Section 3 -- Problem Setting and Threat Model
%
% Job: make the setting precise enough that Section 4 can be formal.
%   (a) the MCP mechanics,          (b) actors and who is protected from whom,
%   (c) how harm arises,            (d) the server's observation set,
%   (e) assumptions and scope,      (f) the decision problem, stated formally.
%
% Budget: ~650 words + one small table.
%
% TERMINOLOGY: "resource" is reserved for r, the thing a call touches. The
% server is the "protected party", never "the protected resource", and the word
% "asset" does not appear here. Sec.4/5 still say "asset register" and
% "(tool, asset) cell" -- one sweep would make the paper uniform.
%
% THE ROUTES-TO-HARM LIST is now ONE sentence with grouped citations: Sec.1 P2
% and Sec.2 P1/P2 already carry the same six keys, and this was the third
% telling of the same list.
%
% The risk-adaptive AC pattern is established here and nowhere else: the old
% Sec.4.7 (Enforcement) that re-stated it with the same three keys has been
% removed, since Sec.4 now stops at rho(c).
% ---------------------------------------------------------------------------
\section{Problem Setting and Threat Model}
\label{sec:threat}

\textbf{The setting.}
An MCP server publishes a catalog of tools over JSON-RPC. A client uses
\texttt{tools/list} to obtain each tool's name, description, and argument
schema, and then invokes a tool with concrete
arguments~\cite{anthropic2024mcp,mcp2026spec}. A tool is only an interface to
an operation. The actual target may be an organizational file, database,
service, or other protected resource. When an agent invokes a tool, the MCP
server performs the requested operation on that resource.


\textbf{Actors.}
The organization deploys the MCP server, owns the connected assets, and defines
the policies used to classify them. The user gives a task to the agent. The
agent chooses a tool and provides the arguments for the call. The MCP server
receives the call, maps it to the relevant asset, and performs the requested
operation. 


\textbf{How harm arises.}
We assume no adversary. Harm can still arise during normal use. The user may
give an incomplete or unclear request. When the agent converts this request into
a tool call, it may select the wrong tool or produce arguments that do not match
the task. Prior work reports parameter errors and task deviations caused by
ambiguous requests and model limitations~\cite{xiong2025butterfly}. An agent
may also understand that an action is risky but still perform it in a concrete
situation~\cite{tang2025riskknowledge}.

Some of these mistakes still produce a schema-valid, authenticated, and
permitted call. Such calls matter because MCP applications can reach system,
network, file, and memory resources~\cite{li2025privilege}. Our work assesses
the possible consequence of the call without trying to determine where the
mistake began. Deliberate attacks by users, agents, or MCP servers are outside
our scope.


\textbf{What the server observes.} At assessment time---after authorization but before execution---the server receives the five quantities in Table~\ref{tab:observables}, together with the caller's authorization context. Identity is not included among the five because the same authorized caller may issue either an appropriate or a mistaken invocation, so identity alone does not distinguish their risk. Moreover, the standard MCP \texttt{tools/call} request carries the tool name and arguments, but no protocol-defined field containing the principal's instruction or the agent's intent~\cite{mcp2026spec}. From the invocation alone, the server therefore cannot determine whether an unusually large transfer is a typo or an approved operation. The estimate below is consequently a function of the five observed quantities, not of what the agent meant.


\begin{table}[t]
  \caption{What an MCP server can observe about an invocation before executing
  it. Intent is not among them.}
  \label{tab:observables}
  \begin{tabular}{@{}cl@{}}
    \toprule
    & Observable \\
    \midrule
    $t$ & the tool named in the request \\
    $r$ & the resource(s) that request resolves to \\
    $a$ & the argument values, valid against the tool's schema \\
    $x$ & session context: the calls already made on this connection \\
    $h$ & the prior invocations already recorded for this deployment \\
    \bottomrule
  \end{tabular}
\end{table}

\textbf{Assumptions.} \emph{(i)} Authentication and transport security are assumed to work correctly. We focus only on the possible consequence of executing the request.
\emph{(ii)} The tool catalog is truthful --- a tool's name,
description, and schema describe what invoking it does --- and no party acts
with intent to harm. A call may be wrong for the task; it is not an attack.
Poisoned descriptions and a compromised server are a separate problem and are
scoped out. \emph{(iii)} The organization provides a classification policy based on security
categorization standards~\cite{nist2004fips199,nist2008sp80060}. The policy
defines resource classes by their possible impact and gives rules for deciding
which class each resource belongs to. The organization registers the resources
used by the MCP server, but does not assign a sensitivity score to each one in
advance. Instead, the policy rules are used to classify both existing and newly
created resources.
\emph{(iv)} The server uses the selected tool and its arguments to identify the
resource affected by the call. \emph{(v)} The server records the calls it
receives, so previous calling history $h$ is available when the risk is
assessed. Our framework only reports the risk score. The organization decides
whether to log it, send it for review, or use it to block a call.

\textbf{The assessment problem.}
We represent an invocation as
$c=\langle t,r,a,x,h\rangle$, because its risk depends on the full call and not
only on the selected tool. We model the result of existing access control as
$\mathrm{perm}(c)\in{0,1}$. Our method considers only permitted calls, for
which $\mathrm{perm}(c)=1$.

The goal is to compute a numerical consequence score $\rho(c)$ that ranks
permitted calls by their possible harm. Our work reports this score but does
not decide whether a call should be allowed, reviewed, or blocked.
Risk-adaptive access control can use a calculated risk value when making such
decisions~\cite{cheng2007fuzzymls,kandala2011radac}, but each organization
decides how to use the score.

The score has two main properties. First, it is numerical rather than limited
to labels such as low, medium, and high. This supports finer comparison between
calls and allows an organization to choose its own thresholds. Second, the
score is explainable: its factors are reported separately so that security
teams can inspect and discuss them. Section~\ref{sec:framework} explains how
$\rho(c)$ is calculated.
% ---------------------------------------------------------------------------
% Section 4 -- The Framework
%
% Budget: 1,400-1,700 words (WRITING-GUIDE.md Sec.5). Largest prose section.
%
% Structure: 4.1 Framework Overview (why the two layers are separate kinds of
% evidence) -> 4.2 The Static Layer (4.2.1 Tool Impact, which absorbs the old
% rules-to-model hand-off; 4.2.2 Resource Sensitivity; 4.2.3 Blast Radius) ->
% 4.3 The Dynamic Layer, ending at rho(c).
%
% Enforcement is out of scope. Sec.3 states that the framework reports the
% score and does not decide whether a call is allowed, reviewed, or blocked, so
% Sec.4 must not instantiate a threshold gate.
%
% Narrative contract inherited from the Introduction:
%   * tension 1 "the inventory does not exist"  -> paid off in Sec.4.2.2
%   * tension 3 "is the model just guessing?"   -> set up in Sec.4.2.1,
%     paid off by the 55/55 result in the Evaluation, NOT here.
%
% Every number below appears verbatim in GROUND-TRUTH.md. Citations are sparse
% by design (WRITING-GUIDE.md Sec.6: method sections cite only what grounds a
% design decision).
%
% Terminology: "resource" throughout, matching Sec.1 and Sec.3 (Table
% tab:observables writes the resource r). The policy's inventory of resources
% is "the register".
%
% NOTE for the next build pass: this section resolves \ref{sec:threat} and
% \ref{sec:eval}, so 03-05 must all be \input in main.tex before the
% cross-references typeset.
% ---------------------------------------------------------------------------
\section{The Framework}
\label{sec:framework}

\framework{} computes the score $\rho(c)$ of \S\ref{sec:threat} in two layers.
A \emph{static} layer estimates, before any particular call occurs, the
inherent consequence of every (tool, resource) pairing a server exposes; a
\emph{dynamic} layer refines that baseline at request time with the concrete
arguments supplied and the deployment's own history of prior calls. The output is an
auditable score and an ordinal band. The framework assesses risk; it does not
define an allow, review, or block policy.

\subsection{Framework Overview}
\label{sec:overview}

The design choice underneath both layers is to keep capability separate from
context. The same tool has very different consequences depending on which
resource it reaches, and the same (tool, resource) pairing becomes more or less
concerning depending on the concrete request and the caller's recent behaviour.
These are different questions, and collapsing them into a single opaque model
decision loses the ability to say which of the two produced a high score.

The static layer holds the facts that are stable across calls: what operation
the tool performs, how sensitive the resource is, and how far one use of that
tool against that resource can propagate. Because none of these depends on a
particular invocation, they are computed ahead of time and reused, giving every
pairing a baseline consequence an operator can inspect factor by factor. The
dynamic layer is separate because the remaining evidence does not exist until a
request arrives: the argument values actually supplied, and whether this action
is usual for the deployment that issued it. Keeping the
two apart also means transient behavioural evidence never rewrites the
underlying classification --- an anomalous call raises the assessment of that
call, not the tier of the tool that served it.

The static layer reads exactly two documents: the tool catalog the server
advertises over \texttt{tools/list}, and the section of the organization's
policy that covers this server. Neither is a \framework{}-specific artifact:
the server already publishes the catalog, and a security team already maintains
the policy. The policy supplies a register of the resources behind the server
and the tools that touch them, named classes, and
the rules by which the organization's own staff recognise them; it states no
severity number anywhere. The organization's per-resource sensitivity table
never reaches the scorer. Here it is an evaluation instrument only --- the
artifact a deployment would not have, which is what lets it serve as an oracle.
\S\ref{sec:eval} separates the real captures from what was written for this
study.

Three primitives carry the static baseline, and they ask deliberately different
questions: impact is about the verb, what one call does; sensitivity about the
noun, what the call touches; blast radius about the quantifier, how far the
consequence travels. Each names a source of consequence the other two cannot
express, which is why none of them is redundant. They are multiplied rather
than averaged (Eq.~\ref{eq:score}), which keeps a low factor influential: a
destructive tool aimed at a worthless resource cannot reach the top of the
range, since only a call severe on every axis gets there. The dynamic layer
then adjusts that product with request-time evidence to produce $\rho(c)$.
\framework{} stops there --- how an organization turns a band into an
operational action is its own decision, not part of the framework.

\subsection{The Static Layer}
\label{sec:static}

The static layer builds a reusable baseline for every (tool, resource) pairing
the server exposes. It estimates three primitives, each on a 1--5 scale, and
multiplies them:
\begin{equation}
\mathrm{score} = \mathrm{sensitivity} \times \mathrm{blast} \times
                 \mathrm{impact}, \qquad \mathrm{score} \in [1, 125].
\label{eq:score}
\end{equation}
Impact is a property of the tool and sensitivity a property of the resource,
while blast radius is a property of the pairing, because how far a consequence
travels depends on both the operation and what it reaches. The layer runs when
the catalog, the register, or the classification policy changes; it is not
recomputed per invocation.

A scan issues four \emph{kinds} of prompt: stage~0 infers the server's domain
once, fixing the vocabulary for what follows; stage~1 assigns tool impact per
tool, rules first and the model only on abstention; stage~2 assigns resource
sensitivity per register row; stage~3 assigns blast radius per (tool, resource)
cell. The last stage dominates a scan's prompt count --- 26 tools over 20
resources is 520 cells. Everything afterwards is deterministic. An assembly pass reconciles
twin tools (a bulk variant of a singular tool, or a deprecated alias of a
canonical one, is tied to its twin so neither can be priced below the other)
and applies a blast floor; a banding function then maps the product onto
$\{\textrm{low}, \textrm{medium}, \textrm{high}, \textrm{critical}\}$. Banding
is not a raw threshold: it encodes explicit security floors --- any
irreversible operation is at least medium --- derived by measuring where an LLM
reviewer systematically disagreed with raw thresholds. Determinism buys exactly
one thing, and we claim only that: band boundaries, twin reconciliation and the
floor do not drift between runs. The residual run-to-run variance sits in the
LLM stages, and blast radius is where it concentrates.

\subsubsection{Tool Impact}
\label{sec:impact}

Tool impact asks one question of a tool's own declaration: is this operation a
read, a write, or a removal? Table~\ref{tab:ladder} is the resulting ladder,
matched by a vocabulary of 281 patterns --- 16 for \emph{none}, 52 metadata, 46
read, 122 write, 45 remove --- plus 119 ambiguous single words matched against
the tool name only, 10 generic read verbs, and 9 parameter detectors. That last
scoping matters: matched across a description instead, the ambiguous words make
a listing tool look destructive.

\begin{table}[t]
\caption{The tool-impact ladder. One question --- read, write, or removal ---
decides the tier.}
\label{tab:ladder}
\small
\begin{tabular}{@{}cp{0.80\columnwidth}@{}}
\toprule
Tier & Operation \\
\midrule
1 & \textbf{none} --- the server talks about itself: ping, health, version,
    whoami \\
2 & \textbf{metadata} --- names, ids, counts, sizes, timestamps, permissions,
    schema, listings; consumption state \\
3 & \textbf{content read or limited write} --- append a line, add a comment,
    set one named field \\
4 & \textbf{ordinary write} --- the caller supplies what the item says \\
5 & \textbf{removal or execution} --- delete, wipe, drop, purge, truncate;
    execute code; move money \\
\bottomrule
\end{tabular}
\end{table}

The tier 3/4 boundary --- limited write against ordinary write --- is not
invented here: two published standards draw it in the same place. HTTP
separates \texttt{PATCH}, a partial modification leaving unmentioned fields
untouched~\cite{dusseault2010patch}, from \texttt{PUT}, where the client sends
the complete representation~\cite{fielding2022http}; CVSS~v4.0 separates
integrity impact \texttt{VI:L}, ``the amount of modification is limited'', from
\texttt{VI:H}, ``a total loss of integrity''~\cite{first2023cvss4}. Writing a
file is a 4; writing one sentence into it is a 3.

Breadth is deliberately absent from the ladder: how many items a call reaches
is coverage, which blast radius scores, so the rules read no breadth vocabulary
(\emph{all}, \emph{every}, \emph{bulk}, \emph{batch}), no glob and no recursion
switch. Parameter detectors are recorded as capability
flags but never move a tier: what a caller \emph{could} pass is not what a call
\emph{does} pass, which is a request-time fact. MCP tool annotations are
recorded the same way, as evidence and never as a
bound~\cite{mcp2026spec} --- an annotation is a component's self-assessment
rather than an independent measurement, and no component sets the ceiling on
its own risk score.

Which scorer decides the tier is settled by a confidence schedule, so the model
enters this axis only where the rules cannot see far enough.
The rules report a confidence with every tier: \textbf{0.35} when no tier verb
matched, \textbf{0.8} when a verb fired, and \textbf{0.95} when a verb fired and
an MCP annotation corroborated it. The model is consulted only below
\textbf{0.5}. An annotation therefore raises confidence, and can suppress the
model's review, but never sets or caps the tier itself.

The model is asked where the rules do not know --- a write whose declaration
never says whether it is bounded --- not where two rules disagree, which the
precedence order settles. What the
arrangement buys is attributability rather than call volume: stages~2 and~3
already spend a prompt per cell, so the tool axis is not where a scan's prompt
volume lives. Every tier the rules decide traces to a named pattern, and a per-tool
field records which scorer decided, so no tier is unattributable. On the
catalogs we scan, the rules answer most of the tool axis and the model is
consulted on the rest; \S\ref{sec:eval} reports the split and where it falls.

\subsubsection{Resource Sensitivity}
\label{sec:sensitivity}

The input a risk model most wants is the one an organization is least able to
publish. A per-resource sensitivity inventory is expensive to build, sensitive
in itself --- filenames can be protected health information --- and stale the
moment a new resource appears. What organizations do publish is a data
classification policy.

\framework{} therefore derives sensitivity by \emph{classify-then-map},
splitting the labour so neither party does the other's job. The organization
states each class's consequence of loss in its own words ---
statutory breach, patient injury, privilege waiver, recoverable embarrassment
--- and writes no numbers. \framework{} owns the 1--5 scale, maps that
adverse-impact language onto it, and classifies each register row into one of
the organization's classes, using the row's description and the policy's own
recognition and aggregation rules. Classify-then-map follows the security
categorization procedure of FIPS~199 and
SP~800-60~\cite{nist2004fips199,nist2008sp80060} --- impact language first,
category second --- applied to a document the organization already maintains
rather than to a labelled inventory. A rule, unlike an inventory row, still
classifies a resource created after the policy was written --- a property of
the construction, not a result we measure.

\subsubsection{Blast Radius}
\label{sec:blast}

Blast radius is scored per (tool, resource) pairing: how far one call's
consequences propagate, which subjects and systems are different afterwards.
The rubric numbers that propagation. (1)~Almost nobody depends on what was
touched; (2)~a few subjects; (3)~a group, the people or systems
attached to what was touched, and the normal case; (4)~the whole population
inside the resource boundary, the subjects the register says one resource
holds; (5)~beyond that boundary. Tier~5 requires a route the organization's own
register sanctions: the resource is a hub other systems
authenticate against, what the call returns is usable on its own elsewhere, one
call reaches the resource's entire population, or the resource is destroyed
with nothing left to restore. Counting items would break on the obvious case:
reading one credential file touches one item and compromises every system that
credential opens.

One deterministic pass then acts on the number: a floor ties blast radius to
the sensitivity of the resource and to irreversible operations. That breaks the
primitives' independence on purpose, in one direction only --- sensitivity can
raise blast, never lower it --- because a call reaching a crown-jewel resource
is not a pinpoint consequence whatever its verb. An earlier version also
applied a roof; it was removed, because a roof can only
ever under-score. Blast radius carries the widest run-to-run variation of the
three stages; \S\ref{sec:limits} gives the scope of the present measurement.

\subsection{The Dynamic Layer}
\label{sec:dynamic}

The static table answers how bad a call would be if it happened; it cannot see
which calls this deployment actually makes, or how large the arguments of this
one are. The dynamic layer supplies those two facts and nothing else: a runtime
signal never gets a scale of its own, it moves a dimension Eq.~\ref{eq:score}
already has. Sensitivity and impact describe the world and one invocation
cannot change them; likelihood is the slot the scan pins at~1, and argument
magnitude is the one respect in which a call may differ from the typical call
the scan priced. Two factors and one bound therefore compose the layer:
\begin{equation}
\rho(c) = \min\Bigl(125,\; \max\bigl(\mathrm{score}(t,r) \times
          \mathrm{lik}(c) \times \mathrm{esc}(c),\;
          \mathrm{floor}(t,r)\bigr)\Bigr),
\label{eq:rho}
\end{equation}
with $\mathrm{lik} \in [0.1, 1.0]$ and $\mathrm{esc} \in \{0.9, 1.0, 1.3\}$.
Only the static score can make $\rho(c)$ large: likelihood can only shrink it,
escalation nudges it either way, and the floor stops the shrinking, so nothing
observed at runtime invents risk the scan did not already price.

\textbf{Likelihood} asks how usual this action is here, as share of traffic
rather than count --- an agent that makes ten thousand calls does not thereby
earn trust on all of them. With $n$ calls recorded at this deployment, $k_t$ of
them on tool $t$, catalog size $V$, and $\pi(t)$ the tool's share of the
catalog,
\begin{equation}
p = \frac{k_{t} + \alpha\,\pi(t)}{n + \alpha}, \qquad
\mathrm{lik} = 0.1 + 0.9\,
  \mathrm{clamp}\!\left(\frac{\log_2 (1/p)}{\log_2 (n+V)},\, 0,\, 1\right).
\label{eq:lik}
\end{equation}
Rarity enters as surprisal because traffic shares span orders of magnitude, and
normalising by the largest surprisal the history admits keeps the value
comparable as that history grows. The prior removes cold start: at $\alpha =
30$ every tool has support from the first call, and a few dozen observations
are enough for evidence to dominate. The key is (deployment, tool) --- not the
resource, which is already inside the cell, and not the individual agent, since
in our corpora no agent spans more than a single day, at a median of 32 calls,
so per-agent history never accumulates. A tool never seen, a deployment too new
to judge, and a call whose resource did not resolve all return $\mathrm{lik} =
1.0$: absence of evidence is not evidence of safety.

\textbf{Escalation} asks whether this call is large \emph{for this deployment}.
A \texttt{limit} of 30 is unremarkable on one server and unheard of on another,
so the deployment's own history sets the cut points, as Tukey's fences $[Q_1 -
1.5\,\mathrm{IQR},\; Q_3 + 1.5\,\mathrm{IQR}]$ over the values passed for that
argument~\cite{tukey1977eda}: above the upper fence 1.3, below the lower 0.9,
between them 1.0. Quartiles are rank-based, so they assume no distribution and
are not dragged around by the very extremes they exist to find, and Tukey's 1.5
leaves roughly 0.7\% of well-behaved values outside, so the bin fires rarely by
construction. Only quantities are binned --- array lengths, and numerics the
scan marked with a critical trigger --- never identifiers: PR~\#25 is not
larger than PR~\#7, and binning it would manufacture risk out of a meaningless
order. Below 20 observations the factor returns 1.0.

\textbf{The floor} is the lower boundary of the band the cell occupied
\emph{before} anything moved it --- an absolute score, and read from the
pre-change cell, so escalation cannot also raise the floor beneath the call it
raises. The runtime may reorder calls inside a band, and may push one up into a
higher band, but can never discount a call out of the band the scan put it in.
The asymmetry is deliberate: wrongly demoting a call lets something through,
while wrongly promoting one costs a reviewer's attention. CVSS~v3.1 bounds its
temporal modifiers the same way, to at most a fifth off the base
score~\cite{first2019cvss31}. So a cell of 27 that a tool accounting for 8\% of
this deployment's traffic discounts to 14.2 is reported as 17, the bottom of
its own band --- moved within the band, not out of it.

$\rho(c)$ is reported with the evidence that produced it, in the organization's
own terms: the three primitives of Eq.~\ref{eq:score}, the policy class the
resource was mapped to (\S\ref{sec:sensitivity}), the operation tier and the
scorer that produced it (\S\ref{sec:impact}), the share of traffic that
discounted the result, and the fence the arguments crossed. Two of the three
primitives are model-assigned, so this is not a model-free audit trail; what
the decomposition gives an operator is a named class and a named operation to
dispute, in place of a single opaque number. Session-level evidence ($x$ in
Table~\ref{tab:observables}) is observable and is not part of this composition,
and \S\ref{sec:limits} gives the scope of the present results. The
framework ends here: it reports risk and prescribes no enforcement mechanism.



% ---------------------------------------------------------------------------
% Section 5 -- Early Results, and Section 6 -- Conclusion.
%
% Source of every number: GROUND-TRUTH.md Sec.9, derived from
%   reports/experiments/v5/five_level_v2_policy_v5r_nacombo/
% (CORPUS_RESULTS.md, live_run.md, live_scale_run.md, <server>_matrix.csv).
% Nothing here is rounded or re-derived.
%
% Style contract for this section: report what the run produced. State facts
% about the setup and the output. Do not editorialize about the quality of the
% work, and do not argue with the numbers -- a reader of an A*-venue paper
% draws those conclusions from the data and the stated scope.
%
% Terminology: "resource" throughout, matching Sec.1, Sec.3 and Sec.4. The
% artifacts call these "assets"; that is the tooling's word, not the paper's.
% ---------------------------------------------------------------------------
\section{Early Results}
\label{sec:eval}

This section reports the risk surface \framework{} produces on real MCP tool
catalogs. The corpus carries no independent per-resource severity annotation,
so the results characterize the output of the framework and the cost of
producing it; \S\ref{sec:limits} states the scope of the setting.

\textbf{Corpus.} We ran one configuration --- resource sensitivity derived from
a single organizational policy document, with no floor coupling the primitives
--- over \textbf{16 server deployments} drawn from \textbf{9 distinct MCP tool
catalogs}. Every catalog is genuine, captured from a running server through the
protocol's own listing call~\cite{mcp2026spec} and used exactly as
returned. The GitHub, Slack, and calendar catalogs are each deployed under two
organizations, and the filesystem catalog under five tenant organizations, with
a different resource register each time. The scan covers \textbf{262 tools},
\textbf{221 register resources}, and \textbf{3{,}615} (tool, resource) cells.
Bands follow the thresholds of \S\ref{sec:framework}: medium at 17, high at 50,
critical at 100, on the 1--125 range of Eq.~\ref{eq:score}.

\begin{table*}[t]
  \centering
  \small
  \caption{The scan, per deployment. \emph{Cells} is tools $\times$ resources;
    \emph{scored} is the subset the relevance gate judged reachable, and the
    band columns partition it.}
  \label{tab:corpus}
  \begin{tabular}{@{}llrrrrrrrrr@{}}
    \toprule
    & & & & & & & \multicolumn{4}{c@{}}{Scored cells by band} \\
    \cmidrule(l){8-11}
    Kind & Deployment & Tools & Res. & Cells & Scored & N/A
         & low & med. & high & crit. \\
    \midrule
    Filesystem   & \texttt{fs\_fintech\_fs}        & 14 & 18 & 252 & 103 & 149 & 39 & 36 & 24 & 4 \\
    Filesystem   & \texttt{fs\_medical\_clinic\_fs}& 14 & 16 & 224 &  86 & 138 & 27 & 37 & 22 & 0 \\
    Filesystem   & \texttt{fs\_corp\_filesystem}   & 14 & 16 & 224 &  87 & 137 & 36 & 26 & 21 & 4 \\
    Filesystem   & \texttt{fs\_law\_firm\_fs}      & 14 & 15 & 210 &  89 & 121 & 33 & 43 & 13 & 0 \\
    Filesystem   & \texttt{fs\_media\_studio\_fs}  & 14 & 16 & 224 & 111 & 113 & 63 & 36 & 12 & 0 \\
    \midrule
    Source repo. & \texttt{github\_real}           & 26 & 20 & 520 & 139 & 381 & 39 & 48 & 52 & 0 \\
    Source repo. & \texttt{github\_helios}         & 26 & 19 & 494 & 126 & 368 & 29 & 46 & 47 & 4 \\
    Chat         & \texttt{slack\_real}            & 16 & 20 & 320 & 101 & 219 & 37 & 48 &  6 & 10 \\
    Chat         & \texttt{slack\_vireo}           & 16 & 15 & 240 &  62 & 178 & 18 & 38 &  5 & 1 \\
    Calendar     & \texttt{calendar\_real}         & 13 & 16 & 208 &  68 & 140 & 22 & 27 & 14 & 5 \\
    Calendar     & \texttt{calendar\_aurora}       & 13 & 17 & 221 &  71 & 150 & 14 & 33 & 18 & 6 \\
    \midrule
    SQL          & \texttt{sqlite\_cbg\_sqlite}    &  5 & 11 &  55 &  37 &  18 & 13 & 13 &  7 & 4 \\
    \midrule
    Public data  & \texttt{finance\_tools}         & 17 &  6 & 102 &  27 &  75 & 18 &  9 &  0 & 0 \\
    Public data  & \texttt{openbb}                 & 30 &  5 & 150 &  36 & 114 & 23 & 12 &  1 & 0 \\
    Public data  & \texttt{sec\_edgar}             & 21 &  6 & 126 &  27 &  99 & 24 &  3 &  0 & 0 \\
    Public data  & \texttt{yahoo\_finance}         &  9 &  5 &  45 &  12 &  33 &  9 &  3 &  0 & 0 \\
    \midrule
    \textbf{Total} & 16 deployments & \textbf{262} & \textbf{221} & \textbf{3{,}615}
      & \textbf{1{,}182} & \textbf{2{,}433} & 444 & 458 & 242 & \textbf{38} \\
    \bottomrule
  \end{tabular}
\end{table*}

\textbf{The risk surface is sparse.} Table~\ref{tab:corpus} gives the result per
deployment. Of the 3{,}615 cells, \textbf{1{,}182 (33\%)} were scored and
\textbf{2{,}433 (67\%)} were returned as not-applicable, meaning the relevance
gate found that the tool cannot reach that resource and no severity is
asserted. The proportion follows the shape of the catalog: the five-tool SQL
server leaves 18 of 55 cells unreachable, while the 30-tool market-data server
leaves 114 of 150, since a catalog of many narrow single-purpose reads spans
few of its resources. Across the scored cells the distribution is 444 low, 458
medium, 242 high, and \textbf{38 critical}. High and critical together are 24\%
of the scored cells and 7.7\% of the full matrix.

\textbf{Most of the tool axis is decided without a model.} The deterministic
ladder decided \textbf{127 of the 151 distinct tools (84\%)}, and the model was
consulted on the remaining 24; counted per deployment, so that a shared catalog
is scored once per organization, the split is 212 of 262 tools by the ladder.
Per catalog, the ladder decided 12 of 14 filesystem tools, 15 of 26 GitHub, 13
of 16 Slack, 9 of 13 calendar, 4 of 5 SQL, and 74 of 77 across the four
public-data servers. The split follows catalog vocabulary: the public-data
catalogs name their operations with generic verbs (get, search, list), whereas
the GitHub catalog uses domain verbs (fork, merge, push) that the read, write,
and remove vocabulary of \S\ref{sec:impact} does not cover, and those are the
tools routed to the model.

\textbf{The top of each matrix.} Table~\ref{tab:top} lists the highest-scoring
cell of every deployment with its three factors. Eight cells in the corpus reach
the maximum of 125. On each of the five filesystem tenants the top cell is
\texttt{write\_file} against the resource that tenant's policy classifies
highest --- the card vault, the security keys, the patient charts, the matter
files, the unreleased imagery --- and the tenant is never named in the scoring
prompt. The three market-data servers that expose only public information
(\texttt{finance\_tools}, \texttt{sec\_edgar}, \texttt{yahoo\_finance}) produce
zero high and zero critical cells across their 273 cells. The single high cell
in that group is \texttt{openbb}'s skill installer writing to the server's own
capability store, at 64; the ladder abstained on that verb and the model scored
it an ordinary write.

\begin{table*}[t]
  \centering
  \small
  \caption{The highest-scoring cell of each deployment, with its
    sensitivity\,$\cdot$\,blast\,$\cdot$\,impact factors. Eight cells across the
    corpus reach the maximum of 125.}
  \label{tab:top}
  \begin{tabular}{@{}llcr@{}}
    \toprule
    Deployment & Highest cell (tool $\times$ resource) & $S{\cdot}B{\cdot}I$ & Score \\
    \midrule
    \texttt{fs\_fintech\_fs}       & \texttt{write\_file} $\times$ \texttt{card-vault}        & 5\,$\cdot$\,5\,$\cdot$\,4 & 100 \\
    \texttt{fs\_corp\_filesystem}  & \texttt{write\_file} $\times$ \texttt{security-keys}     & 5\,$\cdot$\,5\,$\cdot$\,4 & 100 \\
    \texttt{fs\_medical\_clinic\_fs}& \texttt{write\_file} $\times$ \texttt{patient-charts}   & 5\,$\cdot$\,4\,$\cdot$\,4 &  80 \\
    \texttt{fs\_law\_firm\_fs}     & \texttt{write\_file} $\times$ \texttt{matter-files}      & 5\,$\cdot$\,4\,$\cdot$\,4 &  80 \\
    \texttt{fs\_media\_studio\_fs} & \texttt{write\_file} $\times$ \texttt{unreleased-imagery}& 4\,$\cdot$\,4\,$\cdot$\,4 &  64 \\
    \texttt{github\_helios}        & \texttt{push\_files} $\times$ \texttt{helios-scada-gateway} & 5\,$\cdot$\,5\,$\cdot$\,4 & 100 \\
    \texttt{github\_real}          & \texttt{fork\_repository} $\times$ \texttt{infra-config} & 4\,$\cdot$\,5\,$\cdot$\,4 &  80 \\
    \texttt{slack\_real}           & \texttt{conversations\_leave} $\times$ \texttt{incident-response} & 5\,$\cdot$\,5\,$\cdot$\,5 & \textbf{125} \\
    \texttt{slack\_vireo}          & \texttt{usergroups\_users\_update} $\times$ \texttt{usergroup-membership} & 4\,$\cdot$\,5\,$\cdot$\,5 & 100 \\
    \texttt{calendar\_real}        & \texttt{manage-accounts} $\times$ \texttt{connected-account-config} & 5\,$\cdot$\,5\,$\cdot$\,5 & \textbf{125} \\
    \texttt{calendar\_aurora}      & \texttt{delete-event} $\times$ \texttt{aurora-exec}      & 5\,$\cdot$\,5\,$\cdot$\,5 & \textbf{125} \\
    \texttt{sqlite\_cbg\_sqlite}   & \texttt{write\_query} $\times$ \texttt{api\_keys}        & 5\,$\cdot$\,5\,$\cdot$\,5 & \textbf{125} \\
    \texttt{finance\_tools}        & \texttt{calculate} $\times$ \texttt{expression-evaluator}& 4\,$\cdot$\,3\,$\cdot$\,3 &  36 \\
    \texttt{openbb}                & \texttt{install\_skill} $\times$ \texttt{server-capability-install} & 4\,$\cdot$\,4\,$\cdot$\,4 &  64 \\
    \texttt{sec\_edgar}            & \texttt{get\_recommended\_tools} $\times$ \texttt{concept-catalogs} & 3\,$\cdot$\,3\,$\cdot$\,3 &  27 \\
    \texttt{yahoo\_finance}        & \texttt{get\_historical\_stock\_prices} $\times$ \texttt{research-query-pattern} & 3\,$\cdot$\,3\,$\cdot$\,3 &  27 \\
    \bottomrule
  \end{tabular}
\end{table*}

\textbf{Live traffic.} Three of the sixteen deployments --- the GitHub, Slack,
and calendar servers of Table~\ref{tab:live} --- are backed by provisioned
accounts on the vendors' own services, so we also drove them over the protocol. A pilot of 98 live calls (93 \texttt{OK}, 5
\texttt{ERROR}) was followed by the scaled corpus in Table~\ref{tab:live}:
\textbf{3{,}000 live invocations}, 1{,}000 per organization, of which 2{,}998
returned \texttt{OK}; the two errors are transient calendar \emph{operation was
aborted} responses. Every row executed against a real server and a real account
through the same stdio JSON-RPC client the framework uses, with identifiers
discovered from the live estate. Irreversible operations were capped per run at
eight merges, eight repository creations, and five posts to a restricted
channel. The run created and then removed 92 files on protected branches, 80
branches, 39 issues, 5 pull requests, and 8 repositories, and a post-run residue
audit returned zero on every category. Five verbs across the three catalogs
cannot execute on these deployments --- a self-owned fork the vendor's schema
rejects, three Slack \texttt{usergroups\_*} writes that require a paid
workspace, and a calendar RSVP that refuses an organizer --- and are recorded as
errors. The corpus is the runtime input the dynamic layer of
\S\ref{sec:dynamic} consumes; its benign, misuse, and attack labels are assigned
by the generator from each server's policy section.

\begin{table}[t]
  \centering
  \small
  \caption{The scaled live corpus: 3{,}000 real invocations against the three
    provisioned deployments, at a 50/25/25 benign/misuse/attack mix.}
  \label{tab:live}
  \begin{tabular}{@{}lrrrr@{}}
    \toprule
    Deployment & Calls & B / M / A & Tools live & \texttt{OK} \\
    \midrule
    \texttt{github\_helios}   & 1{,}000 & 500 / 250 / 250 & 24 of 26 & 1{,}000 \\
    \texttt{slack\_vireo}     & 1{,}000 & 500 / 250 / 250 & 12 of 16 & 1{,}000 \\
    \texttt{calendar\_aurora} & 1{,}000 & 500 / 250 / 250 & 12 of 13 &    998 \\
    \midrule
    \textbf{Total} & \textbf{3{,}000} & 1{,}500 / 750 / 750 & 48 & \textbf{2{,}998} \\
    \bottomrule
  \end{tabular}
\end{table}

\textbf{Cost.} The static layer spends one prompt per tool, one per register
resource, and one per reachable cell, and the assembly and banding that follow
are deterministic. The dynamic layer's three model-free signals fit in under
0.5\,s per server and score a call in roughly 0.05\,ms. A seventeenth server,
with 119 tools, aborted the inference backend on its first call --- a CUDA
flash-attention kernel abort at that prompt length with a quantized KV cache ---
and is excluded from every number above.

\subsection{Scope of these results}
\label{sec:limits}

The setting bounds what the numbers above cover. The corpus has no independent
per-resource severity table, so Tables~\ref{tab:corpus}--\ref{tab:top} report
the distribution the framework produces rather than agreement with an external
reference, and no baseline scorer or factor ablation is run alongside it. The
2{,}433 not-applicable cells are not independently checked. Run-to-run variance
of the blast-radius stage is not measured in this configuration. All sixteen
registers derive from one classification policy, so behaviour under a different
organization's policy is not covered. The live corpus is reported as captured
traffic: its category labels come from the generator rather than from an
independent annotator, and the optional LLM judge of \S\ref{sec:dynamic} was not
run. Measuring agreement with severities assigned independently by a deploying
organization is the next step.

% ---------------------------------------------------------------------------
% Section 6 -- Conclusion
% ---------------------------------------------------------------------------
\section{Conclusion}
\label{sec:conclusion}

An MCP server authorizes a call before it can know what the call will cost. We
presented \framework{}, which computes that cost in two layers: a static layer
that prices every (tool, resource) pairing from the tool's operation, the
resource's sensitivity, and the blast radius of the two together, and a dynamic
layer that adjusts the baseline with the arguments, the session, and the calling
agent's history. Its distinguishing input is the sensitivity term. Rather than
requiring a labeled inventory that most organizations do not maintain,
\framework{} classifies each register resource against the organization's own
data classification policy and maps the policy's adverse-impact language onto
the 1--5 scale, following the security categorization procedure of FIPS~199 and
SP~800-60~\cite{nist2004fips199,nist2008sp80060}.

Over 16 deployments of 9 real MCP catalogs, the framework priced 3{,}615 cells
from two documents per server, marked two thirds of them unreachable, placed
7.7\% of the full matrix at high or critical, and decided 84\% of the tool axis
with deterministic rules; on the three deployments backed by provisioned vendor
accounts we captured 3{,}000 live invocations to exercise the same servers at
runtime. Every score decomposes into a named policy class, a named operation
tier, and the scorer that produced it, so an organization can dispute a factor
instead of a number. Establishing agreement with severities an organization
assigns independently, and evaluating the dynamic layer against annotated
traffic, are the immediate next steps.


% ------------------------------------------------------------ bibliography ---
% DRAFTING AID -- remove before submission. Forces every curated reference into
% the bibliography, so the full reading set is visible while sections 2+ are
% still unwritten.
\nocite{*}

\balance
\bibliographystyle{ACM-Reference-Format}
\bibliography{refs}

\end{document}
